% !TEX program = xelatex
% ============================================================
%  خطة الظهور في أوائل نتائج جوجل باسم «يَسِّر»
%  (Search Console · السلطة · صفحات الهبوط · العلامة التجارية)
%  الاستخدام على Overleaf:
%    1) New Project ← Blank Project
%    2) الصق المحتوى في main.tex
%    3) Menu ← Compiler ← اختر: XeLaTeX   (مهم للعربية)
%    4) Recompile ← نزّل PDF
% ============================================================
\documentclass[11pt]{article}

\usepackage{polyglossia}
\setmainlanguage{arabic}
\setotherlanguage{english}
\newfontfamily\arabicfont[Script=Arabic,Scale=1.0]{Amiri}
\newfontfamily\arabicfontsf[Script=Arabic,Scale=1.0]{Amiri}
\setmonofont{DejaVu Sans Mono}[Scale=0.8]

\usepackage{geometry}
\geometry{a4paper, margin=2.2cm}

\usepackage{xcolor}
\usepackage{graphicx}
\usepackage{booktabs}
\usepackage{longtable}
\usepackage{array}
\usepackage{enumitem}
\usepackage{titlesec}
\usepackage{fancyhdr}
\usepackage{hyperref}
\usepackage{tcolorbox}
\tcbuselibrary{skins,breakable}
\usepackage{fontawesome5}

% ── الألوان (هوية يَسِّر) ──
\definecolor{gold}{HTML}{B8933B}
\definecolor{teal}{HTML}{0F3D3E}
\definecolor{darkbg}{HTML}{101414}
\definecolor{okgreen}{HTML}{27AE60}
\definecolor{warnred}{HTML}{C0392B}
\definecolor{softgray}{HTML}{F4F4F4}
\definecolor{midgray}{HTML}{7A7A7A}
\definecolor{codebg}{HTML}{1A2020}

\pagestyle{fancy}
\fancyhf{}
\renewcommand{\headrulewidth}{0.4pt}
\renewcommand{\footrulewidth}{0.4pt}
\fancyhead[R]{\small\color{gold}خطة الظهور في أوائل جوجل — منصة يَسِّر}
\fancyhead[L]{\small\color{midgray}Google Ranking Roadmap}
\fancyfoot[C]{\small\thepage}

\titleformat{\section}{\Large\bfseries\color{gold}}{\thesection}{0.6em}{}
\titleformat{\subsection}{\large\bfseries\color{darkbg}}{\thesubsection}{0.5em}{}
\titleformat{\subsubsection}{\normalsize\bfseries\color{midgray}}{\thesubsubsection}{0.4em}{}

\setlist[itemize]{rightmargin=0pt, itemsep=2pt, topsep=3pt}
\setlist[enumerate]{rightmargin=0pt, itemsep=2pt, topsep=3pt}

\hypersetup{
  colorlinks=true, linkcolor=gold, urlcolor=gold,
  pdftitle={خطة الظهور في أوائل جوجل}, pdfauthor={منصة يَسِّر}
}

% ── صناديق مخصّصة ──
\newtcolorbox{goalbox}[1][]{
  colback=gold!8, colframe=gold, boxrule=0.8pt,
  arc=3pt, left=8pt, right=8pt, top=6pt, bottom=6pt,
  breakable, #1
}
\newtcolorbox{notebox}[1][]{
  colback=softgray, colframe=midgray, boxrule=0.6pt,
  arc=3pt, left=8pt, right=8pt, top=6pt, bottom=6pt,
  breakable, #1
}
\newtcolorbox{warnbox}[1][]{
  colback=warnred!6, colframe=warnred, boxrule=0.8pt,
  arc=3pt, left=8pt, right=8pt, top=6pt, bottom=6pt,
  breakable, #1
}
\newtcolorbox{codebox}[1][]{
  colback=codebg, colframe=codebg, coltext=white,
  arc=3pt, left=8pt, right=8pt, top=6pt, bottom=6pt,
  breakable, #1
}

\begin{document}

% ============================================================
\begin{center}
  {\Huge\bfseries\color{gold} خطة الظهور في أوائل جوجل}\\[4pt]
  {\large\color{darkbg} كلمة «يَسِّر» كعلامة تجارية + عبارات محددة}\\[6pt]
  {\small\color{midgray} منصة يَسِّر الفقهية — \today}
\end{center}

\vspace{4pt}
\hrule
\vspace{10pt}

% ============================================================
\section{التشخيص — قبل أي وعد}

\begin{warnbox}
\textbf{حقيقة يجب قولها صريحةً:} كلمة «يسر» عربية شائعة جداً (بمعنى السهولة، وترد في القرآن).
التنافس عليها وحدها على المرتبة الأولى \textbf{غير واقعي عملياً} لمشروع جديد.
\textbf{الهدف الواقعي القابل للتحقيق} هو تصدّر
\textbf{«يَسِّر» كعلامة تجارية} والعبارات المركّبة («منصة يسر» · «يَسِّر الفقه» · «الفقه الميسر»).
\end{warnbox}

\begin{goalbox}
\textbf{الخبر الجيد:} الأساس التقني للـ SEO \textbf{متقدّم بالفعل} في المنصة —
JSON-LD غني (\texttt{WebSite}+\texttt{Organization}+\texttt{Article}+\texttt{Book}+\texttt{FAQPage}+\texttt{BreadcrumbList})،
\texttt{SearchAction}، sitemap بـ\texttt{lastmod}، وأيقونات. الفجوة ليست تقنية، بل
\textbf{السلطة (Authority)} و\textbf{القياس}.
\end{goalbox}

\subsection*{الفجوات المرصودة}
\begin{itemize}
  \item لا تحقّق Google Search Console في \texttt{index.html} — \textbf{لا قياس أصلاً}.
  \item \texttt{sameAs} في \texttt{Organization} فارغ — لا ربط بحسابات اجتماعية.
  \item لا روابط خارجية (Backlinks) — أهم عامل ترتيب مفقود.
  \item عناوين المسائل غير مُحسَّنة لعبارات البحث.
\end{itemize}

% ============================================================
\section{المرحلة ١ — تأسيس القياس (أساس كل شيء)}

\begin{notebox}
بلا قياس، لا يمكن معرفة الترتيب ولا تحسينه. هذه أول خطوة إلزامية.
\end{notebox}

\begin{itemize}
  \item تسجيل الموقع في \textbf{Google Search Console} (نطاق \texttt{fqh.me}).
  \item إضافة \texttt{<meta name="google-site-verification">} في \texttt{index.html}.
  \item إرسال \texttt{sitemap.xml} في GSC.
  \item ربط GSC بحساب GA4 (من خطة اكتمال الموقع).
  \item طلب فهرسة يدوية للصفحة الرئيسية من ``URL Inspection''.
\end{itemize}

% ============================================================
\section{المرحلة ٢ — سلطة النطاق (Off-page)} 

\begin{goalbox}
هذه \textbf{أهم مرحلة} لرفع الترتيب، وهي الغائبة تماماً اليوم. الترتيب في جوجل
يعتمد قبل كل شيء على ``الثقة'' التي تمنحها مواقع أخرى بذكر المنصة.
\end{goalbox}

\begin{itemize}
  \item \textbf{روابط خارجية (Backlinks):} ذكر المنصة في منتديات فقهية، مدوّنات دعوية،
        أدلة المواقع الإسلامية، والمجتمعات العلمية الموثوقة.
  \item \textbf{\texttt{sameAs}:} ربط حسابات المنصة الرسمية في \texttt{Organization} JSON-LD
        (الحقل موجود وبنية القيمة جاهزة — ينتظر الروابط فقط).
  \item \textbf{حسابات اجتماعية رسمية} موثّقة (تخدم الـ \texttt{sameAs} و\textbf{Knowledge Panel}).
  \item \textbf{محتوى قابل للمشاركة:} الفوائد المصوّرة (الميزة الموجودة) تُشارَك فتُنتج إشارات طبيعية.
\end{itemize}

% ============================================================
\section{المرحلة ٣ — صفحات هبوط إستراتيجية}

بدل التنافس على «يسر» العامة، نستهدف عبارات واقعية لها نية بحث واضحة:

\vspace{6pt}
\renewcommand{\arraystretch}{1.35}
\begin{longtable}{@{}p{0.34\textwidth} p{0.58\textwidth}@{}}
\toprule
\textbf{العبارة المستهدفة} & \textbf{الصفحة / الأداة} \\
\midrule
\endhead
أحكام الطهارة · الصلاة · الزكاة · الصيام & صفحات الكتب الثابتة (\texttt{/books/...}) — موجودة ويُحسَّن عنوانها ووصفها \\
منصة يسر · يَسِّر الفقه · الفقه الميسر & الصفحة الرئيسية + \texttt{/about/} + \texttt{alternateName} في JSON-LD \\
حكم [مسألة] · [مسألة] في الفقه & صفحات المسائل الثابتة — إثراء العناوين بعبارات بحثية \\
\bottomrule
\end{longtable}
\renewcommand{\arraystretch}{1.0}

\begin{itemize}
  \item استثمار \texttt{FAQPage} الموجود لجذب \textbf{المقتطفات المميزة} في نتائج جوجل.
  \item تحسين الوصف الميتا (120--160 حرفاً) لكل صفحة.
  \item الروابط الداخلية بين الكتب والمسائل (موجودة) تُوزّع السلطة داخلياً.
\end{itemize}

% ============================================================
\section{المرحلة ٤ — المحتوى الطازج والعميق}

\begin{itemize}
  \item نشر مسائل جديدة دورياً (آلية لوحة الإدارة موجودة).
  \item \texttt{lastmod} في sitemap (موجود) يُسرّع إعادة الزحف عند كل تحديث.
  \item عمق المحتوى: نص المسألة + تعليق الشيخ (موجود) — إشارة جودة.
\end{itemize}

% ============================================================
\section{المرحلة ٥ — Core Web Vitals}

\begin{itemize}
  \item قياس LCP / CLS / INP عبر PageSpeed Insights بعد تركيب GA4.
  \item المنصة قوية هنا مسبقاً (خطوط محلية · PRPL · تجزئة حزم) — الغرض مراقبة عدم التراجع.
  \item تجنّب أي تراجع ناتج عن سكربتات التحليلات (تحميل غير معطِّل).
\end{itemize}

% ============================================================
\section{المرحلة ٦ — إشارات العلامة التجارية «يَسِّر»}

\begin{itemize}
  \item \textbf{Knowledge Panel:} تقوية \texttt{Organization} بـ\texttt{sameAs} + \texttt{logo}
        (الشعار موجود) ← قد تُنشئ جوجل بطاقة معرفة.
  \item \textbf{اسم الموقع في النتائج:} مدعوم بـ\texttt{WebSite.name} (موجود) وأيقونات الموقع (موجودة).
  \item \textbf{\texttt{notranslate}} (موجود) يمنع تشويه الاسم بالترجمة الآلية.
  \item \textbf{تعجيل الفهرسة:} طلب إعادة فهرسة بعد أي تغيير جوهري في الاسم/البنية.
\end{itemize}

% ============================================================
\section{المرحلة ٧ — القياس الأسبوعي}

\begin{itemize}
  \item GSC ← ``الأداء'': تتبّع ظهور كلمة «يَسِّر» والعبارات المستهدفة أسبوعياً.
  \item مراقبة ظهور \textbf{Knowledge Panel} تدريجياً.
  \item فحص دوري بـ\textbf{Rich Results Test} لكل نوع schema.
  \item مراجعة \textbf{Core Web Vitals} شهرياً في GSC.
\end{itemize}

% ============================================================
\section{الجدول الزمني والأثر}

\renewcommand{\arraystretch}{1.35}
\begin{longtable}{@{}p{0.08\textwidth} p{0.44\textwidth} p{0.14\textwidth} p{0.24\textwidth}@{}}
\toprule
\textbf{\#} & \textbf{المرحلة} & \textbf{المدة} & \textbf{زمن الأثر} \\
\midrule
\endhead
١ & تأسيس القياس (GSC) & نصف يوم & فوري \\
٢ & سلطة النطاق (Off-page) & مستمر & 4--12 أسبوعاً \\
٣ & صفحات هبوط إستراتيجية & يومان & 2--6 أسابيع \\
٤ & المحتوى الطازج & مستمر & متراكم \\
٥ & Core Web Vitals & نصف يوم & فوري \\
٦ & إشارات العلامة & يوم & 4--8 أسابيع \\
٧ & القياس الأسبوعي & أسبوعي & — \\
\bottomrule
\end{longtable}
\renewcommand{\arraystretch}{1.0}

\begin{warnbox}
\textbf{لا يمكن تسريع الفهرسة والتصنيف قسراً.} الفهرسة تحتاج \textbf{٢--٦ أسابيع}،
وبناء السلطة يحتاج \textbf{أشهراً}. أي وعد بغير ذلك غير صادق.
\end{warnbox}

% ============================================================
\section{معايير القبول}

\begin{itemize}
  \item[${\color{okgreen}\faCheckCircle}$] المنصة موثّقة في Google Search Console ومفهرسة.
  \item[${\color{okgreen}\faCheckCircle}$] \texttt{sameAs} مملوء بحسابات المنصة الرسمية.
  \item[${\color{okgreen}\faCheckCircle}$] صفحات الكتب والمسائل لها عناوين وأوصاف مُحسَّنة.
  \item[${\color{okgreen}\faCheckCircle}$] «يَسِّر» تظهر كاسم الموقع في النتائج.
  \item[${\color{okgreen}\faCheckCircle}$] ظهور معرفي (Knowledge Panel) عند تراكم الإشارات.
  \item[${\color{okgreen}\faCheckCircle}$] \texttt{check:seo} وRich Results Test بلا أخطاء.
\end{itemize}

\end{document}
